% GERADO por scripts/analyse_feedback.py — NÃO EDITAR À MÃO.
% Correr o script outra vez depois de a recolha fechar; o texto acompanha os dados.
\subsection{Utilidade percebida dos alertas entregues}
\label{sec:av_feedback}

O canal passou a acompanhar cada alerta entregue com dois botões, \emph{útil} e \emph{não ajudou}, e a registar o voto de quem carrega. As regras de análise foram fixadas antes de existir um único voto e não foram alteradas depois: mínimo de vinte votos efetivos para reportar qualquer proporção, um voto por pessoa e por alerta com o último a substituir o anterior, salvaguarda que repete o cálculo sem o votante dominante quando este representa mais de quarenta por cento dos votos, sempre sujeita ao mesmo mínimo, e intervalos de confiança de Wilson, apropriados a proporções com amostras pequenas.
Contaram apenas votos sobre alertas presentes no registo partilhado.

Foram registados 31 votos válidos entre 2026-09-01 e 2026-09-03, que correspondem a 20 votos efetivos de 2 pessoas distintas, sobre 16 alertas. 5 votos foram posteriormente alterados pela mesma pessoa, e apenas o último de cada par conta. 4 votos foram excluídos por não corresponderem a nenhum alerta do registo partilhado.

Dos 20 votos efetivos, 19 classificaram o alerta como útil, ou seja 95\%, com intervalo de confiança de Wilson a 95\% entre 76\% e 99\%. A largura deste intervalo é a medida honesta do que 20 votos permitem afirmar.

Uma única pessoa representa 80\% dos votos efetivos, acima do limite de 40\% fixado no protocolo. Excluindo-a, restam 4 votos, dos quais 4 classificam o alerta como útil.
Este recorte também fica abaixo do mínimo de 20; nenhuma segunda proporção é reportada e a amostra não sustenta uma leitura independente do votante dominante.

Quatro limitações acompanham este resultado e nenhuma delas se resolve com mais votos. A primeira é a autosseleção: vota quem quer, e quem considera um alerta indiferente tende a não carregar em nada, o que empurra a amostra para os extremos. A segunda é a ausência de contrafactual, uma vez que nenhum grupo recebeu a variação de preço sem a explicação que a acompanha; nada aqui atribui a utilidade à explicação em si. A terceira é a distância entre utilidade percebida e decisão melhor, que é a hipótese fundadora deste trabalho e permanece por testar. A quarta é o desconhecimento de quem vota: o canal é público e não se sabe se quem responde pertence ao público que a dissertação assume.
